\documentclass[11pt]{article}
\usepackage{amsmath, amssymb}
\usepackage[margin=1in]{geometry}
\usepackage{booktabs}
\usepackage{graphicx}
\PassOptionsToPackage{obeyspaces,spaces}{url}
\usepackage{url}
\usepackage[hidelinks]{hyperref}
\setlength{\parskip}{4pt}
\setlength{\parindent}{0pt}
\graphicspath{{model_output/}{model_output/moist_v2/}}

\newcommand{\py}{\partial_y}
\newcommand{\Shat}{\hat S}
\newcommand{\Hhat}{\hat H}
\newcommand{\Wq}{W_q}
\newcommand{\Wstar}{W^*}
\newcommand{\vhalf}{v_{\rm half}}
\newcommand{\vslab}{v_{\delta}}

\title{Calibrating the moist Sobel--Schneider model against ERA5:\\
what the reanalysis settles, and what it does not}
\date{2026-08-01}

\begin{document}
\maketitle

\section*{Summary}

The model's moisture, thermodynamic and momentum equations all contain a
velocity written $v$, and they do not mean the same velocity by it. The
moisture equation's $-(2a{-}1)vW$ needs two branches that between them exhaust
the column, so each holds half the column mass. The thermodynamic equation's
$(d\Delta_z/H)\py v$ describes a layer of depth $d$ at the top of the
troposphere and another at the bottom, with $v=0$ in between. Those two
velocities differ by the ratio of their depths, a factor of about 2.5 at the
observed surface pressure, and every coefficient defined as ``transport per
unit $v$'' scales with that factor. Measuring one coefficient under one reading
and pairing it with another measured under the other reading is what produced
the negative gross moist stability that has been the obstacle since 2026-07-31.
It was an accounting error, not a physical result.

Two combinations of the coefficients do not depend on the reading at all:
the ratio $\Hhat/\Shat$, and the column
$\Wstar = \Shat/(L_v(2a{-}1))$ at which the gross moist stability changes sign.
Those are the right things to compare with observations, and they can be
compared without ever deciding what $v$ means.
ERA5 gives $\Wstar$ in the range 44 to 57~kg\,m$^{-2}$, depending on how the
time mean is taken, against 44.3~kg\,m$^{-2}$ for the model at its default
$a=0.85$. The model's default therefore sits at the extreme
strong-transport end of what the observations allow, and it is the one choice
that puts a $W_c = 50$ run below its own sign change, since such a run settles
at a quiescent column of 50.7~kg\,m$^{-2}$.

Acting on that: a V2 run at $W_c=50$ with $a$ lowered to 0.77 was integrated to
convergence for this report. The aggregated regime of 2026-07-31, five
precipitating bands, a 78~mm\,day$^{-1}$ equatorial spike and a bone-dry third
of the domain, does not appear. The run instead reproduces the recognizable
single-ITCZ circulation of the $W_c=40$ control almost exactly
(Fig.~\ref{fig:acheck}).

The model's gross dry stability, which I reported on 2026-07-31 as 2.55 times
too small, is not. That claim came from dividing the observed dry-static-energy
flux by a half-column velocity while comparing against a layer-velocity
$\Shat$. Read consistently, the model's $\Shat = Cd\Delta_z/H = 7.75\times
10^{7}$~J\,m$^{-2}$ corresponds to a slab pair 195~hPa deep, which the observed
$[v]$ profile accommodates comfortably. Measured against the branch depth the
observed wind shape itself implies, 238~hPa, the model's $\Shat$ is 18\% low.
Both that claim and the 2026-07-31 recommendation of $a=0.95$ are withdrawn.

What ERA5 does not settle is the sign of the model's gross moist stability at
its current defaults. Four defensible ways of taking the time mean give
$a$ between 0.77 and 0.84 and $\Wstar$ between 44 and 57~kg\,m$^{-2}$, a spread
that straddles the quiescent column. Section~\ref{sec:aggregation} lays that
out; Section~\ref{sec:open} lists what would narrow it.

\section{What the model asks of the observations}
\label{sec:ask}

Three terms in the model transport something meridionally, and all three are
written in terms of one velocity $v$ (Appendix~\ref{app:eqs} has the full
equations):
%
\begin{align}
\text{moisture:} &\quad \py\!\left[-(2a-1)\,v\,W\right],
\label{eq:moist}\\
\text{thermodynamic:} &\quad \frac{d\,\Delta_z}{H}\,\py v,
\label{eq:thermo}\\
\text{momentum:} &\quad \text{the factor of 2 in the $v$ equation.}
\label{eq:mom}
\end{align}

Equation~\eqref{eq:moist} is built from a two-branch column. A fraction $a$ of
the column water vapor $W$ sits in the lower branch and moves at $-v$; the
remaining fraction $1-a$ sits in the upper branch and moves at $+v$. The
column-integrated flux is then
$(1-a)W(+v) + aW(-v) = -(2a-1)vW$. For the fractions to be $a$ and $1-a$ the
two branches must exhaust the column between them. Since they also have equal
and opposite velocities, and time-mean mass balance requires their mass
transports to be equal and opposite, they must have equal masses. Two branches
of equal mass that exhaust the column each hold half of it, so the interface is
at $p_s/2$ and the branch velocity is
%
\begin{equation}
\vhalf = \frac{2g\,\Psi_{\rm ext}}{p_s},
\label{eq:vhalf}
\end{equation}
%
with $\Psi_{\rm ext}$ the observed overturning mass transport, that is, the
zonal-mean mass streamfunction evaluated at the level where $[v]$ reverses
sign.

Equation~\eqref{eq:thermo} is built from a different column. It comes from
vertical advection $w\,\partial_z\theta$ with $w = d\,\py v$, the vertical
velocity that continuity gives when the horizontal flow occupies a layer of
depth $d$, and $\partial_z\theta = \Delta_z/H$. That picture has a moving layer
of depth $d$ at the top, another at the bottom, and $v=0$ between them. A slab
of pressure depth $dp$ carrying mass transport $\Psi_{\rm ext}$ moves at
%
\begin{equation}
\vslab = \frac{g\,\Psi_{\rm ext}}{dp}.
\label{eq:vslab}
\end{equation}

Dividing~\eqref{eq:vslab} by~\eqref{eq:vhalf},
%
\begin{equation}
\frac{\vslab}{\vhalf} = \frac{p_s/2}{dp}.
\label{eq:ratio}
\end{equation}
%
At the observed tropical $p_s = 997$~hPa and a slab depth near 195~hPa this
ratio is 2.55. Any coefficient defined as a transport divided by $v$ inherits
it exactly. That single factor is the whole of the discrepancy I reported on
2026-07-31.

\section{The part of the answer that does not depend on the reading}
\label{sec:invariant}

Before choosing between the two readings, it is worth isolating what does not
require the choice. The gross moist stability is
%
\begin{equation}
\Hhat = \Shat - L_v(2a-1)W,
\label{eq:hhat}
\end{equation}
%
and both $\Shat$ and $(2a-1)$ are transports per unit $v$, so both scale with
the same factor~\eqref{eq:ratio}. Two combinations therefore do not scale at
all:
%
\begin{equation}
\Wstar \;\equiv\; \frac{\Shat}{L_v(2a-1)}
      \;=\; \frac{\Shat\,W}{\Shat - \Hhat},
\qquad
\frac{\Hhat}{\Shat} \;=\; 1 - \frac{W}{\Wstar}.
\label{eq:invariants}
\end{equation}
%
$\Wstar$ is the column at which $\Hhat$ changes sign. In V1 and V2 the
equilibrium column is nearly uniform at the quiescent value
$\Wq = W_c + \tau_c E_0$ set by the local evaporation--precipitation balance,
so the comparison of $\Wq$ with $\Wstar$ alone decides which regime a run is
in. The second form of $\Wstar$ in~\eqref{eq:invariants} needs only the two
regressed stabilities and never names a value of $2a-1$, so it cannot import a
normalization along with it.

Measured over $|\varphi|\le 20^\circ$ from ERA5 2000--2019, with the barotropic
mass adjustment applied to $[v]$ (Section~\ref{sec:stability}):
%
\begin{equation}
\Wstar = 56.8~\mathrm{kg\,m^{-2}},
\qquad
\frac{\Hhat}{\Shat} = 0.303 \quad\text{at } W = 39.6~\mathrm{kg\,m^{-2}}.
\end{equation}

Table~\ref{tab:invariant} puts the model beside that. At the model's fixed
$\Shat$, its $a$ is the only free dial in $\Wstar$, and the default $a=0.85$
puts the sign change at 44.3~kg\,m$^{-2}$, below the 50.7~kg\,m$^{-2}$
quiescent column that the V1 defaults produce. Every V2 run at $W_c=50$ was
therefore in the negative-stability corner before it started.

\begin{table}[htbp]
\centering
\caption{The two normalization-free comparisons. $\Hhat/\Shat$ is evaluated at
the quiescent column $\Wq = 50.7$~kg\,m$^{-2}$ that the V1 defaults produce,
which is the value the model actually sits at. The ERA5 row uses the
transport-weighted regression over every month; Section~\ref{sec:aggregation}
gives the spread across other time means.}
\label{tab:invariant}
\begin{tabular}{lcccc}
\toprule
& $a$ & $2a-1$ & $\Wstar$ (kg\,m$^{-2}$) & $\Hhat/\Shat$ at $\Wq$ \\
\midrule
model default        & 0.850 & 0.700 & 44.3 & $-0.145$ \\
ERA5-consistent      & 0.773 & 0.546 & 56.8 & $+0.107$ \\
\bottomrule
\end{tabular}
\end{table}

\section{The $d/H$ argument, laid out}
\label{sec:dh}

This is the step that was unclear in chat, so it is written out slowly here.

\textbf{Step 1. The thermodynamic equation fixes the dry-static-energy flux; it
is not a separate assumption.} The model's temperature equation contains
$(d\Delta_z/H)\py v$. Multiplying by the column heat capacity
$C = c_p \Pi\, \Delta p/g$ turns that term into the divergence of a column
dry-static-energy flux, so
%
\begin{equation}
F_{\rm DSE} = C\,\frac{d\,\Delta_z}{H}\,v \equiv \Shat v,
\qquad
\Shat = C\,\frac{d\,\Delta_z}{H}.
\label{eq:shatmodel}
\end{equation}
%
With $C = 5.164\times10^{6}$~J\,m$^{-2}$\,K$^{-1}$, $d = 4$~km,
$\Delta_z = 60$~K and $H = 16$~km, $\Shat = 7.746\times10^{7}$~J\,m$^{-2}$.
Nothing has been assumed beyond the equation itself.

\textbf{Step 2. $d\Delta_z/H = 15$~K is a temperature difference across one
layer of depth $d$.} That number looks small next to $\Delta_z = 60$~K, and
that is what made the model's $\Shat$ look wrong. But it is only half the
story, because $\Shat$ also contains the full-column $C$.

\textbf{Step 3. Whether 15~K is right depends entirely on what $v$ means.} Write
the same flux as a two-slab transport: mass transport $\Psi_{\rm ext}$ times
the dry-static-energy difference $\Delta s$ between the branches. With
$\Psi_{\rm ext} = v\,dp/g$,
%
\begin{equation}
F_{\rm DSE} = \frac{v\,dp}{g}\,\Delta s
\quad\Longrightarrow\quad
\Shat = \frac{dp\;\Delta s}{g}
\quad\Longrightarrow\quad
\Delta s = \frac{g\,\Shat}{dp}.
\label{eq:dsofdp}
\end{equation}
%
The model's $\Shat$ is one number; the branch temperature difference it implies
is whatever~\eqref{eq:dsofdp} gives for the slab depth in question. Equating
$\Delta s = c_p \Pi\, \Delta\Theta$ and using $\Delta p = p_s - p_t$,
$\Delta\Theta = (\Delta p/dp)(d\Delta_z/H)$. At $dp = \Delta p/2$, that is 30~K.
At $dp = 195$~hPa, it is 62~K, comparable to $\Delta_z$ itself. The
``15~K'' is a per-layer number that only becomes a branch contrast once a layer
depth is named.

\textbf{Step 4. The observed $\Delta s$ settles it.} ERA5's dry-static-energy
flux per unit overturning is $F_{\rm DSE}/\Psi_{\rm ext} = 38.9$~kJ\,kg$^{-1}$
over $|\varphi|\le20^\circ$. Putting that into~\eqref{eq:dsofdp} with the
model's $\Shat$ gives $dp = 195$~hPa. So the model's gross dry stability is
exactly the statement that its two branches are slabs 195~hPa deep,
39~kJ\,kg$^{-1}$ apart in $s$.

\textbf{Step 5. The alternative I floated in chat was wrong.} I suggested that
the thermodynamic equation should carry $(H-d)\Delta_z/H = 45$~K instead of
15~K. That reproduces the literal-bulk half-column $\Shat$ to 2\%
($2.324\times10^{8}$ against $2.283\times10^{8}$), which is why it looked
right. It is the wrong fix, because it repairs a layer-velocity equation by
matching a half-column measurement. The two-readings table of the next section
is the correct picture.

\section{The two readings, measured}
\label{sec:readings}

\begin{figure}[htbp]
\centering
\includegraphics[width=\textwidth]{era5_normalization_coefficients.png}
\caption{Every transport coefficient against the slab depth that defines $v$.
Purple, blue and green are ERA5. Red dashed in the left panel is the model's
$\Shat = Cd\Delta_z/H$, a fixed number, whose intersection with the observed
curve at 195~hPa is the depth the model's dry stability corresponds to. Green
vertical line: the branch depth the observed $[v]$ profile implies on its own,
238~hPa. Grey dotted: the half-column reading, $dp = p_s/2$. Right panel, red
dashed: what $\Hhat$ becomes if the model's $\Shat$ is paired with a $2a-1$
measured at some other depth, which is the mixing that produced the negative
gross moist stability; the black dot at 195~hPa is the only self-consistent
pairing. Centre panel, orange: what $2a-1$ would be for two genuinely uniform
slabs of that depth, against the blue coefficient the observed flux actually
demands.}
\label{fig:coeff}
\end{figure}

Table~\ref{tab:readings} is Fig.~\ref{fig:coeff} in numbers. Every column is
linear in $dp$, as~\eqref{eq:ratio} requires, and $\Wstar$ is constant down the
table, as~\eqref{eq:invariants} requires. The half-column row is the $dp =
499$~hPa member of the same family, and its $2a-1 = 1.394$ exceeds the ceiling
$2a-1 \le 1$ that follows from $a \le 1$: a signal, visible at the time, that
the half-column reading was not the model's.

\begin{table}[htbp]
\centering
\caption{Transport coefficients under each reading of $v$. ERA5 2000--2019,
$|\varphi|\le20^\circ$, mass-adjusted $[v]$, regressed over all months. Observed
tropical $p_s = 997$~hPa and $W = 39.6$~kg\,m$^{-2}$.}
\label{tab:readings}
\begin{tabular}{lccccc}
\toprule
reading & $\Shat$ (J\,m$^{-2}$) & $2a-1$ & $a$ & $\Hhat$ (J\,m$^{-2}$)
& $\Wstar$ (kg\,m$^{-2}$) \\
\midrule
half-column, $dp = p_s/2$ & $1.978\times10^{8}$ & 1.394 & 1.197 & $5.99\times10^{7}$ & 56.7 \\
slab, $dp = 150$~hPa      & $5.943\times10^{7}$ & 0.419 & 0.709 & $1.80\times10^{7}$ & 56.8 \\
slab, $dp = 196$~hPa      & $7.766\times10^{7}$ & 0.547 & 0.774 & $2.35\times10^{7}$ & 56.8 \\
slab, $dp = 250$~hPa      & $9.906\times10^{7}$ & 0.698 & 0.849 & $3.00\times10^{7}$ & 56.8 \\
slab, $dp = 300$~hPa      & $1.189\times10^{8}$ & 0.838 & 0.919 & $3.60\times10^{7}$ & 56.8 \\
\bottomrule
\end{tabular}
\end{table}

Mixing the readings is what produced the negative gross moist stability.
Holding the model's layer-velocity $\Shat$ fixed and pairing it with each
candidate coefficient:

\begin{center}
\begin{tabular}{lr}
\toprule
combination & $\Hhat$ (J\,m$^{-2}$)\\
\midrule
model $\Shat$ with the half-column flux-matched $2a-1 = 1.394$ & $-6.07\times10^{7}$\\
model $\Shat$ with the half-column mass partition $2a-1 = 0.900$ & $-1.17\times10^{7}$\\
model $\Shat$ with the slab $2a-1 = 0.546$ at the matched depth & $+2.34\times10^{7}$\\
\bottomrule
\end{tabular}
\end{center}

\section{Equal pressure depths, and how thick the branches really are}
\label{sec:depth}

Spencer's correction of 2026-08-01 was that the two layers must have equal
depths in \emph{pressure}, not equal geometric depths, and it is what makes the
slab picture self-consistent. Equal $dp$ gives equal mass, and equal mass is
what equal and opposite branch velocities require. Measured in ERA5, a 4~km
layer spans 153~hPa just below the model's tropopause and 367~hPa just above
the ground, a mass ratio of 2.4, so the model's own geometry read literally
cannot give its two layers equal mass. Fig.~\ref{fig:vert} panel (a) draws both
readings on the observed $[v]$ profile.

\begin{figure}[htbp]
\centering
\includegraphics[width=\textwidth]{era5_vertical_structure.png}
\caption{(a) The observed $[v]$ at $\pm12^\circ$, composited into one cell,
with the model's two layers drawn two ways: equal geometric depth of 4~km each
(green dashed, 153 and 367~hPa, unequal mass) and equal pressure depth of
195~hPa each (purple, the depth the model's $\Shat$ implies). (b) $\Psi$
vanishes at both ends once the column mass flux is removed, which is the
statement that the two branch transports are equal and opposite; the level
where $[v]$ changes sign is what divides them. (c) The observed $s(p)$, whose
branch contrast the transport must reproduce.}
\label{fig:vert}
\end{figure}

Three routes lead to a slab depth, and they do not agree:

\begin{center}
\begin{tabular}{lr}
\toprule
route & depth\\
\midrule
from the observed $s(p)$ profile, matching $\Delta s = 38.9$~kJ\,kg$^{-1}$ & 90~hPa\\
from the model's own $\Shat$, through~\eqref{eq:dsofdp} & 195~hPa\\
from the shape of the observed $[v]$, with no energy data used & 238~hPa\\
\bottomrule
\end{tabular}
\end{center}

The third is the one that makes the comparison a test rather than a fit. For
each branch of the annual-mean $[v]$, split at the level of nondivergence, the
participation ratio $\left(\int|v|\,dp\right)^2 / \int v^2\,dp$ returns the
depth of the top hat that carries the same transport at the same
peak-equivalent speed, and it uses only the wind. Averaged over
$5^\circ$--$25^\circ$ it gives 296~hPa for the upper branch, 180~hPa for the
lower, 238~hPa in the mean. Against the observed $\Shat$ at that depth,
$9.44\times10^{7}$~J\,m$^{-2}$, the model's $\Shat$ is 18\% low.

The first route disagreeing with the others is the honest residual of the
two-slab idealization. Two 195~hPa slabs taken from the observed $s(p)$ profile
differ by 33.7~kJ\,kg$^{-1}$, 13\% short of the 38.9 the transport needs,
because the real branches correlate $[v]$ with $s$ inside themselves and a
uniform slab cannot. That is a real limitation on how precisely $d$ can be
pinned, and it is a 13\% effect rather than the factor of 2.5 that mixing
normalizations produced.

One consequence for branch speed: the observed peak annual-mean overturning
carries 0.51~m\,s$^{-1}$ as a half-column branch velocity, 1.31~m\,s$^{-1}$ as
a 195~hPa slab, and 1.07~m\,s$^{-1}$ as a 238~hPa slab. The model's $v$ is a
slab velocity, so model-versus-ERA5 comparisons of $\max|v|$ belong against the
last two. The 0.49~m\,s$^{-1}$ quoted in my 2026-07-31 notes was the
half-column number and should not have been compared with the model.

\section{The moisture coefficient}
\label{sec:a}

There are three different things one can mean by $a$, and they answer three
different questions.

\textbf{The mass partition} is $a = W_{\rm lower}/W_{\rm total}$ at a chosen
interface. It answers ``where is the water''. At the half-mass level $p_s/2$ it
is 0.9501 over $|\varphi|\le20^\circ$ (ERA5 1979--2023), and it barely moves:
0.9477 within $10^\circ$, 0.9506 within $30^\circ$, 0.9501 globally. Its
seasonal cycle in the tropical mean spans 0.0012, from 0.9496 in March to
0.9508 in October, because the hemispheres cancel; at a fixed latitude the
swing is twenty times larger, 0.939 to 0.964 at $15^\circ$N with the maximum in
local winter. Across 45 annual means it spans 0.9466 to 0.9529 with a trend of
$-0.0011$ per decade. Every one of those variations is dwarfed by the choice of
interface: 0.982 at 400~hPa, 0.950 at 500, 0.883 at 600, 0.772 at 700.

\textbf{The flux-matched coefficient} is the regression slope of the observed
column moisture flux on the model's transport basis $vW$. It answers ``what
coefficient reproduces the observed transport'', which is what the model needs,
and it absorbs the correlation between $[v]$ and $[q]$ within each branch. It
depends on the reading of $v$, and that is the whole content of
Table~\ref{tab:readings}.

\textbf{The literal two-slab partition} is
$(W_{\rm lower\;slab} - W_{\rm upper\;slab})/W$ for slabs of depth $dp$
anchored at the ground and at the model's tropopause. It answers ``what would
two genuinely uniform slabs carry''. At $dp = 196$~hPa it is 0.588 against a
flux-matched 0.547, so the uniform-slab idealization would over-transport by
7\%. The sign of that correction is opposite to the half-column case, where the
flux-matched 1.401 exceeded the literal 0.900 by 56\%, because the half-column
lower branch is 500~hPa deep and its water sits at the very bottom where $|v|$
is largest, whereas a 196~hPa slab is thinner than the real lower branch and
reaches less far into dry air.

\begin{figure}[htbp]
\centering
\includegraphics[width=\textwidth]{era5_normalization_partition.png}
\caption{The literal two-slab coefficient at $dp = 195$~hPa across latitude and
season (left), and its annual mean against the flux-matched value and the
model's default (right). The seasonal structure is the monsoon signature: the
lower slab holds a larger share of the column in the winter subtropics, where
subsidence and the trade inversion trap water low.}
\label{fig:partition}
\end{figure}

\begin{figure}[htbp]
\centering
\includegraphics[width=\textwidth]{era5_a_calibration_transport.png}
\caption{Annual-mean column moisture transport, with $v$ read as the model
reads it. Left: the observed flux against the model's structure at the
flux-matched coefficient and at the spec value $a = 0.85$. Right: the
coefficient the observed flux demands, now comfortably below its own ceiling.
Under the half-column reading this curve sat at 1.6 to 1.8, above the ceiling,
which was the visible symptom of the mismatch.}
\label{fig:transport}
\end{figure}

\begin{figure}[htbp]
\centering
\includegraphics[width=\textwidth]{era5_a_calibration_vertical.png}
\caption{Why $a$ depends so strongly on the interface: the deep-tropical $[q]$
and $|[v]|$ profiles and the cumulative column water. Half the column mass sits
above 499~hPa; about 5\% of the column water does.}
\label{fig:vertical}
\end{figure}

\section{How the time mean is taken, and why it is the largest uncertainty}
\label{sec:aggregation}

The regressions above pool every month and latitude, so they weight by
transport and are dominated by the strong monsoon months. The model is steady
and has no seasonal cycle, so a time-mean circulation and its time-mean
transport are the closer counterpart. Four defensible choices give
Table~\ref{tab:aggregation}.

\begin{table}[htbp]
\centering
\caption{The moisture coefficient and the regime boundary under four ways of
taking the time mean, all at $dp = 195$~hPa, ERA5 2000--2019,
$|\varphi|\le20^\circ$. $\Hhat/\Shat$ is at the observed
$W = 39.6$~kg\,m$^{-2}$.}
\label{tab:aggregation}
\begin{tabular}{lccccc}
\toprule
time mean & $\Shat$ (J\,m$^{-2}$) & $2a-1$ & $a$ & $\Wstar$ (kg\,m$^{-2}$)
& $\Hhat/\Shat$\\
\midrule
every month, transport-weighted & $7.746\times10^{7}$ & 0.546 & 0.773 & 56.8 & 0.303\\
annual mean, $\bar v\,\bar W$   & $7.581\times10^{7}$ & 0.589 & 0.795 & 51.5 & 0.267\\
equinox months, $\bar v\,\bar W$& $7.579\times10^{7}$ & 0.589 & 0.794 & 51.5 & 0.259\\
annual mean, $\overline{vW}$    & $7.581\times10^{7}$ & 0.683 & 0.842 & 44.4 & 0.107\\
\bottomrule
\end{tabular}
\end{table}

The spread is seasonal covariance. The last row keeps the correlation, within
the seasonal cycle, between a strong lower branch and a moist column, so its
transport basis is the largest and the coefficient it demands is the smallest
per unit transport [plausible]. The second row removes that correlation by
multiplying the two annual means. The equinox row, restricted to March, April,
September and October, is the closest analogue of the model, which is steady,
hemispherically symmetric and has no solstitial cell to average away. That it
lands on the annual-mean value to three decimal places is reassuring: the
worry that the annual-mean tropical overturning is a small residual of two
large opposing cells, and therefore a bad target, does not bite here.

The consequence is that the observations do not settle the sign of the model's
gross moist stability at $W_c = 50$. $\Wstar$ from 44.4 to 56.8~kg\,m$^{-2}$
straddles the 50.7~kg\,m$^{-2}$ quiescent column. What the observations do
settle is comparative, and that is enough to act on: $a = 0.85$ lies at or just
beyond the strong-transport end of the range, and gives the lowest $\Wstar$ of
any of them.

Recommendation: take $a$ from the equinox row, $a \approx 0.79$, and carry
0.77 to 0.84 as its uncertainty.

\section{The gross stabilities and their error budget}
\label{sec:stability}

\begin{figure}[htbp]
\centering
\includegraphics[width=\textwidth]{era5_stability_regression.png}
\caption{The regressions the stabilities come from, per unit half-column branch
velocity. (a) and (b) are the scatters themselves, coloured by latitude, so a
poor fit or a latitude-dependent slope would be visible rather than hidden
inside one number. (c) is the same quantity resolved in latitude instead of
aggregated. The local $\Hhat$ goes negative within a few degrees of the
equator; see the caveat in the text.}
\label{fig:stab}
\end{figure}

Both stabilities are regression slopes against the same velocity, so neither
requires choosing a branch interface. Over four latitude bands
(ERA5 2000--2019, half-column normalization, mass-adjusted):

\begin{center}
\begin{tabular}{lccccc}
\toprule
band & $\Shat_{\rm eff}$ & $\Shat_{\rm bulk}$ & $\Hhat_{\rm eff}$
& $W$ (kg\,m$^{-2}$) & $\Hhat/\Shat$\\
\midrule
$|\varphi|\le10^\circ$ & $1.909\times10^{8}$ & $2.287\times10^{8}$ & $4.93\times10^{7}$ & 44.5 & 0.258\\
$|\varphi|\le15^\circ$ & $1.946\times10^{8}$ & $2.283\times10^{8}$ & $5.46\times10^{7}$ & 42.2 & 0.281\\
$|\varphi|\le20^\circ$ & $1.978\times10^{8}$ & $2.283\times10^{8}$ & $5.99\times10^{7}$ & 39.6 & 0.303\\
$|\varphi|\le30^\circ$ & $1.997\times10^{8}$ & $2.301\times10^{8}$ & $6.35\times10^{7}$ & 35.1 & 0.318\\
\bottomrule
\end{tabular}
\end{center}

Two caveats carry into everything above.

\textbf{The barotropic mass adjustment is the dominant uncertainty on
$\Shat$.} ERA5's zonal-mean $[v]$ on pressure levels carries a spurious net
column mass flux, which leaves $|\Psi(p_s)|$ averaging 232 and peaking at
780~kg\,m$^{-1}$\,s$^{-1}$. Removing the mass-weighted column mean of $[v]$
first changes $\Shat$ by $-33.8\%$, from $2.989\times10^{8}$ to
$1.978\times10^{8}$, while the peak $|\Psi_{\rm ext}|$ moves by only $-10.2\%$
and the peak $|v|$ by $-10.3\%$. $\Shat$ is the most sensitive of the three
because the transported quantity $s$ has a mean near 330~kJ\,kg$^{-1}$, so a
small spurious barotropic $[v]$ fabricates a large spurious flux, while $\Psi$
and $v$ feel only the mass flux itself. Treat 34\% as the floor on the
uncertainty in $\Shat$ until the mass-consistent ERA5 product is used.
Moisture is far more tolerant of the same error because $q$ has a small mean,
which is why the $D$ route in Section~\ref{sec:other} is robust and the energy
route is not.

\textbf{The local $\Hhat$ is not uniformly positive.} Panel (c) of
Fig.~\ref{fig:stab} shows $\Hhat$ falling through zero within about
$8^\circ$ of the equator, reaching roughly $-50$~MJ\,m$^{-2}$ just south of it,
and rising to $+140$~MJ\,m$^{-2}$ by $25^\circ$N. The band aggregate is
dominated by the large-$|v|$ subtropical points. The equatorial values are
poorly conditioned, since $v$ passes through zero there and the local ratio is
a ratio of two small numbers, so do not quote $-50$ as a firm number. This is
the observed counterpart of the slightly negative gross moist stability that
Spencer said might still be of interest.

\section{Does the ERA5-consistent $a$ remove the aggregated regime?}
\label{sec:acheck}

The prediction is direct: hold everything else fixed, lower $a$ from 0.85 to
the ERA5-consistent end, and a $W_c = 50$ run should stop being
negative-stability. One run tests it. The control is the $W_c = 40$ run of
2026-07-31, which reaches positive stability the other way, by lowering the
column instead of raising the crossover.

\begin{figure}[htbp]
\centering
\includegraphics[width=\textwidth]{era5_a_check.png}
\caption{Equilibrium profiles at $ny = 801$, $dt = 30$~s, converged under the
slow-drift gate. Red: the 2026-07-31 configuration, $a = 0.85$ and $W_c = 50$.
Green: the same run with $a = 0.77$. Yellow: the $W_c = 40$ control at
$a = 0.85$. Green and yellow are close in every field; red is the aggregated
regime.}
\label{fig:acheck}
\end{figure}

\begin{table}[htbp]
\centering
\caption{The V2 test. $\Wstar$ and $\Wq$ are the model's own values.
``Bands'' counts contiguous regions with $P > 0.5$~mm\,day$^{-1}$; ``dry'' is
the fraction of the domain where $P$ is identically zero. Mean and eddy MSE
fluxes are peak values in MW\,m$^{-1}$.}
\label{tab:acheck}
\small
\begin{tabular}{lcccccccc}
\toprule
run & $\Wstar$ & $\Wq$ & $\min\Hhat$ & bands & dry & $P_{\max}$
& mean flux & eddy flux\\
& \multicolumn{2}{c}{(kg\,m$^{-2}$)} & (MJ\,m$^{-2}$) & & &
(mm\,day$^{-1}$) & \multicolumn{2}{c}{(MW\,m$^{-1}$)}\\
\midrule
$a=0.85$, $W_c=50$ & 44.2 & 50.66 & $-32.8$ & 5 & 0.36 & 77.8 & $+104.7$ & 106.2\\
$a=0.77$, $W_c=50$ & 57.4 & 50.66 & $+8.3$  & 1 & 0.00 & 7.4  & $+17.0$  & 1.0\\
$a=0.85$, $W_c=40$ & 44.2 & 40.66 & $+4.8$  & 1 & 0.00 & 9.0  & $+14.5$  & 1.5\\
\bottomrule
\end{tabular}
\end{table}

The prediction holds. At $a = 0.77$ the five precipitating bands collapse to
one, the dry gap closes, the peak precipitation falls from 78 to
7.4~mm\,day$^{-1}$, the jet loses its extra easterly notch and weakens from
73.8 to 58.7~m\,s$^{-1}$, and the mean MSE flux returns to a modest poleward
$+17$~MW\,m$^{-1}$ with the eddy flux a 1.0~MW\,m$^{-1}$ bit player, against
$+104.7$ and 106.2 in near-cancellation before. Every one of those numbers
lands next to the $W_c = 40$ control rather than next to the run it was
changed from. The aggregated regime was the negative gross moist stability, and
the negative gross moist stability was the calibration error.

That said, the resulting solution is still not a good match to the observed
circulation. Its jet peaks at 58.7~m\,s$^{-1}$ against an observed 200~hPa
zonal-mean maximum of 31.0~m\,s$^{-1}$, and its meridional temperature contrast
is 1.4 times the observed one at $16^\circ$ and $24^\circ$, 1.6 times at
$33^\circ$
(Table~\ref{tab:score}). Fixing the moisture coefficient removes a pathology;
it does not make V2 realistic, and the overshoot of the jet is common to every
V2 configuration run so far.

\section{The other parameters}
\label{sec:other}

\begin{figure}[htbp]
\centering
\includegraphics[width=\textwidth]{era5_moisture_budget.png}
\caption{(a) The column moisture budget decomposed into its total, mean and
eddy parts. (b) The eddy flux against the gradient it should follow, with the
regression slope, the model default and the postdoc's lat-lon estimate.
(c) The boundary-layer $\theta_e$ contrast against the observed
free-tropospheric $\theta$ contrast; the former is a contaminated proxy for the
radiative-convective target, so it bounds rather than determines it.}
\label{fig:budget}
\end{figure}

\textbf{Eddy diffusivity $D$.} Zonal-mean fields cannot form the eddy moisture
flux directly, because zonal averaging has already destroyed the covariance.
The steady column budget routes around that: the total flux follows from
$E - P$ by meridional integration, and subtracting the mean flux
$\int [v][q]\,dp/g$ leaves the eddy flux, stationary and transient together.
Then $D = -F_{\rm eddy}/(\py W)$. Over ERA5 2000--2009 the median over
$|\varphi|\le25^\circ$ is $7.2\times10^{5}$~m$^2$\,s$^{-1}$ with quartiles
$4.1\times10^{5}$ to $1.8\times10^{6}$, but the latitude structure is the real
result: $D$ rises from $2.7\times10^{5}$ at $5^\circ$S and
$5$--$6\times10^{5}$ near $10^\circ$ to $2.0\times10^{6}$ at $20^\circ$S,
$2.8\times10^{6}$ at $20^\circ$N and $4.4\times10^{6}$ at $25^\circ$N. An order
of magnitude separates the deep
tropics from the subtropics. That reconciles the model default $D = 10^{6}$,
the earlier aggregate near $7\times10^{5}$, and the postdoc's lat-lon tropical
estimate of 1 to $2\times10^{5}$: the model's default is a subtropical value
applied to the deep tropics, and the residual gap to the lat-lon estimate is
plausibly the stationary eddies that the zonal-mean route lumps into the same
term. The budget closes: the total flux reaches only 1.7 and
$-0.5$~kg\,m$^{-1}$\,s$^{-1}$ at $85^\circ$N and S, where it must vanish, once
the global mean $E-P$ imbalance is removed first.

\textbf{Evaporation $E_0$.} ERA5 gives $4.58\times10^{-5}$ within $10^\circ$,
$4.79\times10^{-5}$ within $20^\circ$ and $4.62\times10^{-5}$ within
$30^\circ$~kg\,m$^{-2}$\,s$^{-1}$. The model default is
$4.6\times10^{-5}$. No change is needed.

\textbf{The radiative-convective temperature contrast.} Spencer's position of
2026-08-01, which this report adopts: the relaxation target in V2 should be
radiative-convective equilibrium rather than radiative equilibrium, and he does
not know how he convinced himself otherwise when the spec was written. He is
also right that the target is not observable. A column-by-column RCE free
troposphere sits on the moist adiabat through its own boundary layer, so its
meridional structure would follow the boundary-layer $\theta_e$, whose
equator-to-$24^\circ$ contrast is 16.6~K against 3.2~K for the observed
free troposphere, a factor of 5.2. But boundary-layer $\theta_e$ in a
reanalysis already carries the circulation's influence, and a true
latitude-by-latitude RCE has time-mean $v = w = 0$ by definition, so it is not
a state of the real atmosphere at all. Those two numbers therefore bound
$\Delta\theta^{\rm rad}$ rather than determine it. Treat it as a tuning
parameter with ERA5 setting the limits.

\textbf{$\beta$ and the latitude mapping.} The model's $\beta = 2\times
10^{-11}$~m$^{-1}$\,s$^{-1}$ equals Earth's $2\Omega\cos\varphi/a$ at
$\varphi = 29.1^\circ$. Following Spencer's suggestion to compute a proper
Hadley-range average, $\langle\beta\rangle = 2\Omega\int\cos^2\varphi\,
d\varphi \,/\, \left(a\int\cos\varphi\,d\varphi\right)$ over $0$ to $30^\circ$
gives $2.190\times10^{-11}$, 9.5\% above the model's value; the equatorial
value is $2.289\times10^{-11}$ and the value at $15^\circ$ is
$2.211\times10^{-11}$. For turning model $y$ into a latitude, the Coriolis
mapping $f = \beta y = 2\Omega\sin\varphi$ is the one to use, since it preserves
the dynamically relevant quantity: $y = 2$, 3 and 4~Mm map to $16^\circ$,
$24^\circ$ and $33^\circ$. Arc length on Earth's radius gives $13^\circ$,
$27^\circ$ and $54^\circ$ instead, and the difference at large $y$ is
substantial.

\section{Model against ERA5}
\label{sec:score}

\begin{table}[htbp]
\centering
\caption{Equilibrium model states against ERA5. Temperature contrasts are
$\theta(0)-\theta(y)$ in potential temperature, evaluated at the
Coriolis-mapped latitudes. Peak $|v|$ for ERA5 is quoted as a slab velocity at
$dp = 195$ and 238~hPa, the two depths of Section~\ref{sec:depth}, because that
is the quantity the model's $v$ is; the half-column value is
0.51~m\,s$^{-1}$.}
\label{tab:score}
\begin{tabular}{lccccc}
\toprule
 & $\Delta\theta$ at $16^\circ$ & at $24^\circ$ & at $33^\circ$
 & jet (m\,s$^{-1}$) & peak $|v|$ (m\,s$^{-1}$)\\
\midrule
ERA5                        & 0.56 & 3.21 & 8.75  & 31.0 & 1.07 to 1.31\\
V2, $a=0.85$, $W_c=50$      & 1.69 & 7.69 & 21.33 & 73.8 & 5.86\\
V2, $a=0.77$, $W_c=50$      & 0.76 & 4.49 & 14.20 & 58.7 & 1.88\\
V2, $a=0.85$, $W_c=40$      & 0.90 & 5.13 & 15.72 & 60.8 & 2.27\\
\bottomrule
\end{tabular}
\end{table}

The ERA5-consistent run is the closest of the three to the observations in
every column, and it is still too steep by a factor of 1.4 to 1.6 in the
temperature contrast and too strong by a factor of 1.9 in the jet. The jet
overshoot is shared by every V2 configuration and predates this calibration: it
appeared when the radiative retarget steepened the relaxation target, and it is
a separate problem from the moisture coefficient.

The peak $|v|$ comparison in Table~\ref{tab:score} replaces the one in my
2026-07-31 notes, which mixed conventions by comparing the model's slab
velocity against ERA5's half-column velocity.

\section{Outstanding questions}
\label{sec:open}

\textbf{Things ERA5 could settle but has not, for want of data.} Five ERA5
directories that would remove the largest uncertainties here are present but
empty: \path|col_wv_flux_north|, \path|col_mse_flux_north_st_eddy|,
\path|col_dse_flux_north_st_eddy|, \path|div_col_wv_flux| and
\path|tend_col_wv|. The first three would give the column fluxes directly from
ERA5's own mass-consistent, submonthly accounting, instead of reconstructing
them from monthly zonal means. That would retire the 34\% mass-adjustment
uncertainty on $\Shat$ (Section~\ref{sec:stability}), separate the stationary
from the transient eddies in $D$, and remove the aggregation ambiguity of
Section~\ref{sec:aggregation}, which is currently the single largest unknown in
the calibration. Downloading them is the highest-value next step by a wide
margin.

\textbf{Which time mean is the right calibration target.} Section
\ref{sec:aggregation} argues for the equinox product-of-means on the grounds
that the model is steady and symmetric, and the equinox and annual-mean routes
agree. The mean-of-products route, which keeps the seasonal covariance, gives a
coefficient 16\% larger and puts $\Wstar$ exactly at the model's default value.
Deciding this properly needs a statement of what the steady model is meant to
represent, which is a modeling question rather than an observational one.

\textbf{Parameters still to calibrate.} $\Delta_z$ and $H$ directly from ERA5,
by taking $\theta$ at the lapse-rate tropopause minus $\theta$ at the surface
and the tropopause height, which would check whether 60~K and 16~km hold up;
$\tau$ from the clear-sky radiative cooling fields \path|rlntcs|, \path|rsntcs|
and \path|olr|, all present on disk; $W_c$ and $\tau_c$ from regressing $P$ on
$W$, noting that monthly zonal means give only a smoothed pickup curve and so
bound rather than fit the pair; $y_1$ and $y_0$; and optionally $v_d$ from the
zonal-momentum-budget residual. $\theta_{00}$ is cosmetic
in V2 and needs no calibration.

\textbf{Whether $d$ should be re-specified in pressure.} The model carries $d$
in metres and combines it with $H$ in metres. Everything in
Section~\ref{sec:depth} says the meaningful parameter is the layer's pressure
depth. The two are interchangeable at fixed $\Shat$, so this is a question of
whether to make the model's own statement of its assumptions honest rather than
a question of changing any result.

\textbf{The V2 jet overshoot.} The three configurations of
Section~\ref{sec:acheck} produce jets 1.9 to 2.4 times the observed 200~hPa
maximum, and the $\Delta\theta^{\rm rad}$ ladder of 2026-07-30 reaches higher
still. In that same ladder, latent heating steepened the model's meridional
temperature gradient rather than flattening it, which is the opposite of the
rationale given in the spec; that comparison has not been re-run since. The
overshoot is now the largest remaining model-observation discrepancy, and it is
untouched by this calibration.

\textbf{The slightly negative gross moist stability case.} Spencer flagged
this as possibly still of interest. The observations support it locally:
$\Hhat$ is negative within roughly $8^\circ$ of the equator in ERA5
(Section~\ref{sec:stability}), even though the tropical aggregate is firmly
positive at $\Hhat/\Shat = 0.30$. A model configuration with $\Wq$ just above
$\Wstar$ would be the analogue, and it is now locatable in closed form.

\appendix

\section{Equations and parameter values}
\label{app:eqs}

The V2 equation set, with $\theta_{\rm rad}$ the radiative relaxation target
and $\Lambda = L_v/C$:
%
\begin{align}
\partial_t u + v\,\py u + u\,\py v &= \beta y v - \epsilon_u u
  - \py\!\left(v_d\,H(u)\,u\right),\\
2\,\partial_t v &= -\beta y u - \py \phi + \kappa_v \partial_y^2 v,\\
\partial_t \theta + \frac{d\,\Delta_z}{H}\,H(\py v)\,\py v
  &= \frac{\theta_{\rm rad} - \theta}{\tau} + \Lambda P,\\
\partial_t W + \py\!\left[-(2a-1)vW - D\,\py W\right] &= E_0 - P,
  \qquad P = \frac{(W - W_c)^+}{\tau_c}.
\end{align}

Column constants (\path|ss09/moist_constants.py|):
$C = c_p \Pi\,\Delta p/g = 5.164\times10^{6}$~J\,m$^{-2}$\,K$^{-1}$,
$L_v = 2.5\times10^{6}$~J\,kg$^{-1}$,
$\Lambda = L_v/C = 0.4843$~K per kg\,m$^{-2}$\,s$^{-1}$,
$\Shat = Cd\Delta_z/H = 7.746\times10^{7}$~J\,m$^{-2}$.
The identity $C\Lambda = L_v$ is what makes precipitation cancel from the
column MSE budget.

Parameter values used above: $d = 4$~km, $\Delta_z = 60$~K, $H = 16$~km,
$\beta = 2\times10^{-11}$~m$^{-1}$\,s$^{-1}$, $\tau = 3.197\times10^{6}$~s,
$v_d = 2.5$~m\,s$^{-1}$, $a = 0.85$, $D = 10^{6}$~m$^2$\,s$^{-1}$,
$W_c = 50$~kg\,m$^{-2}$, $\tau_c = 14400$~s,
$E_0 = 4.6\times10^{-5}$~kg\,m$^{-2}$\,s$^{-1}$,
$\Delta\theta^{\rm rad} = 75$~K, $y_1 = 9439$~km. Quiescent column
$\Wq = W_c + \tau_c E_0 = 50.66$~kg\,m$^{-2}$.

\section{Parameter status}
\label{app:status}

\begin{center}\small
\begin{tabular}{ll}
\toprule
status & parameters\\
\midrule
calibrated here & $\Shat$, $\Hhat$, $a$ and $2a-1$, $E_0$, $D$, $\Wstar$\\
calibrated with a stated caveat & $\Delta\theta$ (contaminated proxy), $\beta$
  (a choice of reference latitude)\\
calibratable, not yet done & $\Delta_z$, $H$, $\tau$, $y_1$, $y_0$,
  $(W_c, \tau_c)$, $v_d$\\
over-determined & $d$: two slab parameters against three observables,
  13\% residual\\
not observable & $\Delta\theta^{\rm rad}$ (needs offline RCE), $\theta_{00}$
  (cosmetic)\\
purely numerical & $\epsilon_u$ (regularization), $\kappa_v$\\
\bottomrule
\end{tabular}
\end{center}

The organizing principle behind this table: calibrate the combinations the
equations use, not the individual symbols. Calibrating $a$ on its own is what
pushed the gross moist stability negative, because the symbols only mean
something once the normalization of $v$ is fixed. $\Wstar$ and $\Hhat/\Shat$
are the combinations that survive that ambiguity, and they are where the
comparison with observations belongs.

\section{Data and scripts}
\label{app:data}

ERA5 monthly zonal means at \path|~/Dropbox/data/ecmwf/era5/monthly/|,
symlinked in the repository as \path|era5-data/|. Fields used: \path|hus|,
\path|cwv|, \path|ps|, \path|va|, \path|ua|, \path|ta|, \path|zg|, \path|evap|,
\path|pr|. Precipitation is the only latitude-longitude file (4~GB) and is
streamed month by month, since dask is not installed in this environment.
Year ranges differ by script for cost reasons and are stated with each result:
1979--2023 for the mass partition, 2000--2019 for the stabilities and the
normalization, 2000--2009 for the moisture budget.

Scripts, all run as \path|python scripts/<name>.py| from the repository root:

\begin{center}
\begin{tabular}{ll}
\toprule
script & what it produces\\
\midrule
\path|era5_normalization.py| & Sections~\ref{sec:invariant},
  \ref{sec:readings}, \ref{sec:aggregation}; Figs.~\ref{fig:coeff},
  \ref{fig:partition}\\
\path|era5_vertical_structure.py| & Section~\ref{sec:depth};
  Fig.~\ref{fig:vert}\\
\path|era5_stability_calibration.py| & Section~\ref{sec:stability};
  Fig.~\ref{fig:stab}\\
\path|era5_a_calibration.py| & the mass partition of Section~\ref{sec:a};
  writes the cache\\
\path|era5_a_calibration_figures.py| & Figs.~\ref{fig:transport},
  \ref{fig:vertical}\\
\path|era5_moisture_budget.py| & Section~\ref{sec:other};
  Fig.~\ref{fig:budget}\\
\path|era5_calibration_v2_check.py| & Section~\ref{sec:acheck};
  Fig.~\ref{fig:acheck}, Table~\ref{tab:score}\\
\bottomrule
\end{tabular}
\end{center}

The V2 test run of Section~\ref{sec:acheck} is
\path|model_output/moist_v2/wc50_a077/|, produced by

\begin{quote}\small
\path|run-sw-model --stop-at-steady-state --slow-drift-gate|\\
\path|--ny 801 --dt 30 --backend numba|\\
\path|--enable-moisture --enable-latent-heating|\\
\path|--cwv-frac 0.77 --w-crit 50|
\end{quote}

converging at day 1392, alongside the existing
\path|model_output/moist_v2/m4_v2_dyr75/| and
\path|model_output/moist_v2/wc40/|.

\section{Two verification checks worth recording}
\label{app:checks}

The identity behind~\eqref{eq:moist} is checked numerically: the model's
$-(2a-1)vW$ and the two-layer bulk flux evaluated at the half-mass interface
agree to $2.8\times10^{-14}$~kg\,m$^{-1}$\,s$^{-1}$, which is
$5\times10^{-14}$\% of the peak flux. The two expressions are the same
algebra, and the code agrees.

The pressure-level column integral of $[q]$ runs 0.56~kg\,m$^{-2}$ below
ERA5's own total column water vapour in the tropics, a mean relative error of
1.95\%, with the largest discrepancies over Tibet and the Sahara where zonal
asymmetry in surface pressure is largest. Piecewise-constant against
trapezoidal integration changes $a$ by 0.0006, and removing the below-ground
mask entirely changes it by 0.0007. Spencer's call on 2026-08-01 was that an
error of this size is not a concern.

\end{document}
